% !TeX root = ../../../main.tex
\subsection{Related Work}
\label{ch08:sec:related_work}


\textbf{Clinical temporal information extraction.}
Clinical timeline reconstruction builds on a long tradition of temporal reasoning and temporal information extraction in medicine \cite{ch08:kohane1987temporal,zhou2007temporal}.
Early annotation schemes and shared tasks, including i2b2 and THYME, focused on identifying clinical events, time expressions, and temporal relations such as before, after, overlap, and containment \cite{ch08:uzuner20112010,sun2013evaluating,sun2013annotating,ch08:styler2014temporal}.
Systems such as MedTime and cTAKES-temporal operationalized these ideas for clinical notes by extracting time expressions and temporal relations \cite{lin2013medtime,savova2010mayo,lin2016multilayered}.
Pre-LLM work also studied patient timelines and longitudinal record summarization, including timeline generation from clinical notes and HARVEST for summarizing longitudinal patient records \cite{nikfarjam2013towards,hirsch2015harvest}.
These studies established the importance of temporal structure in clinical narratives, but primarily focused on temporal annotation, relation extraction, or summarization rather than grounding narrative-derived absolute timelines in structured EHR evidence.


\noindent
\textbf{LLM-based clinical timeline reconstruction.}
Recent work has moved from relation-centered temporal extraction toward absolute event timelines.
Leeuwenberg and Moens studied absolute event timeline extraction from English clinical reports, while Frattallone-Llado et al. showed that aligning discharge summaries with structured EHR data can improve inpatient event timeline precision \cite{leeuwenberg2020towards,ch08:frattallone2024using}.
Textual time-series work has further shown that clinical narratives can be converted into event--time tuples for temporal forecasting, but these approaches are primarily text-only and do not use structured EHR rows as calibration evidence during reconstruction \cite{ch08:wang2025large,noroozizadeh2026reconstructing,ch08:noroozizadeh2026forecasting,ch08:kumar2026temporally}.
Recent shared tasks such as ChemoTimelines have renewed interest in patient-level clinical timeline extraction using prompting, fine-tuning, and dictionary-enhanced LLM systems \cite{yao2025chemotimelines,zhang2025uwbionlp}.
Our work differs from these efforts by treating timeline reconstruction as multimodal temporal grounding: narrative text supplies event semantics, while structured EHR rows are retrieved as partial temporal evidence for timestamp calibration.


\noindent
\textbf{Asymmetric evidence in structured and narrative EHR data.}
Clinical machine learning studies have shown that structured and unstructured EHR data encode complementary information for phenotyping, prediction, and patient representation \cite{ch08:moldwin2021empirical,liu2022multimodal,seinen2025using}.
Structured data provide dense timestamped measurements, but often omit symptoms, functional status, patient-reported histories, and clinical reasoning.
Narrative notes preserve this context, but are often retrospective and not organized as chronological trajectories.
This asymmetry makes clinical timeline reconstruction different from standard multimodal fusion: the goal is not to merge two complete views of the same event sequence, but to use structured rows as partial temporal evidence for a narrative-derived patient trajectory.
Our retrieval-augmented scaffolded workflow is designed for this setting, selecting candidate structured rows for each text-derived event and updating timestamps only when the retrieved evidence is clinically relevant.


